\clearpage
\begingroup
\thispagestyle{plain}

% 영문 논문요약 페이지입니다.
% 이 파일의 역할: 영문 논문요약을 작성합니다.
% 국문 논문요약의 내용을 바탕으로 연구 목적, 연구 방법, 주요 결과, 결론을 영어로 옮깁니다.
% 단순 직역보다 학위논문 논문요약에 맞는 명확하고 간결한 문장으로 다듬는 것이 좋습니다.
% 영문 제목, 저자명, 전공, 학교명, 지도교수명은 sub/0-preamble.tex의
% \thesistitleEng, \authorEng, \majorEng, \universityEng, \advisorEng 값을 사용합니다.
% 이 파일에서는 영문 논문요약 본문, Keywords, 하단 제출 문구를 주로 수정합니다.
\begin{center}
\setstretch{1}
\vspace*{10mm}

% 영문 논문요약 제목: 고정 문구입니다.
{\fontsize{14pt}{22.4pt}\selectfont \textbf{A\hspace{0.55em}B\hspace{0.55em}S\hspace{0.55em}T\hspace{0.55em}R\hspace{0.55em}A\hspace{0.55em}C\hspace{0.55em}T}\par}

\vspace{10mm}
% 영문 논문 제목: sub/0-preamble.tex의 \thesistitleEng 값을 사용합니다.
% 줄간격 133%(표지와 동일 비율 28/21): 14pt x 4/3 = 18.67pt.
{\fontsize{14pt}{18.67pt}\selectfont \thesistitleEng\par}

\vspace{10mm}
% 영문 저자명: sub/0-preamble.tex의 \authorEng 값을 사용합니다. (이탤릭체)
{\fontsize{13pt}{20.8pt}\selectfont \textit{\authorEng}\par}

\vspace{10mm}
% 영문 전공명과 학교명: sub/0-preamble.tex의 \majorEng, \universityEng 값을 사용합니다.
{\fontsize{13pt}{20.8pt}\selectfont \majorEng\par}
{\fontsize{13pt}{20.8pt}\selectfont \universityEng\par}
{\fontsize{13pt}{20.8pt}\selectfont \universityLocationEng\par}

\vspace{15mm}
% 영문 지도교수명: sub/0-preamble.tex의 \advisorEng 값을 사용합니다.
{\fontsize{13pt}{20.8pt}\selectfont Supervised by Professor \advisorEng\par}
\vspace{15mm}
\end{center}

\begingroup
\setstretch{1.8}
\normalfont\normalsize\selectfont

% 영문 논문요약 본문: 아래 문단들을 논문 내용에 맞게 수정합니다.
\noindent
The abstract is best written after the full thesis has been completed.
It should present the purpose, method, major findings, and conclusion of the study in a concise form.
The first sentence should clearly state the research topic and purpose, and may also indicate the nature of the method, such as an experiment, survey, theoretical analysis, or simulation.

The method should then be summarized briefly.
If the research design, data collection, or analysis method distinguishes the study from previous work, those features should be emphasized.
When the method itself is not the main contribution of the thesis, it should be described only to the extent needed for readers to understand the findings.

The most important parts of the abstract are the findings and conclusion.
The findings should be reported objectively, without unnecessary inference or overinterpretation.
The conclusion should briefly state how the study answers its research purpose.
Because readers often turn to the abstract to grasp the main findings and significance of the study, the findings and conclusion should occupy a sufficient portion of the abstract.

\vskip 0.5cm
\noindent
% Keywords: 영문 주요어를 쉼표로 구분하여 입력합니다.
\textbf{Keywords:} \thesiskeywordsEng

\vfill

% 하단 제출 문구: 첨부 예시처럼 당구장 표시(※)로 시작합니다.
% 학위수여 연월은 sub/0-preamble.tex의 \ThesisConferralYear, \ThesisConferralMonth에서 수정합니다.
\noindent{\setstretch{1}\fontsize{10pt}{16pt}\selectfont ※ A thesis submitted to the Committee of the Graduate School of Korea National University of Education in partial fulfillment of the requirements for the degree of \degreeEng\space in \ThesisConferralDateEng.\par}
\endgroup
\endgroup
